\section{Path degrees, ranks and uniform estimates}\label{sec:paths}

We prove Theorem~\ref{path:stable}. Retain the notation
$X_j^{(b_0,b_j)}$ from Section~\ref{sec:bundle-families}.
The proof is computer assisted. We first express the divisor degrees by products
of fixed rational matrices and the relevant subspace ranks by a recurrence with
three states. A contraction estimate then reduces all path lengths to a specified
finite calculation. The matrices, recurrence, rational bounds, and error estimate
are given below.

\subsection{Divisor degrees along the path}
Write $a_i=\langle m_i,r_i\rangle\in[-1,1]$ for the coordinate on the
$i$th dual hexagon and $d\nu(a)=(2-|a|)\,da$. The factors in the volume formula are
\[
 L_0=b_0+1-a_0,\qquad L_c=3+a_{c-1}-a_c\ (1\le c<j),
 \qquad L_j=b_j+1+a_{j-1}.
\]
Set
\begin{equation}\label{path:kernel}
 \mathcal K(a,a')=\frac{(3+a-a')^2}{2},\qquad
 \ell_0(a)=\frac{(b_0+1-a)^{b_0}}{b_0!},\qquad
 h_0(a)=\frac{(b_j+1+a)^{b_j}}{b_j!}.
\end{equation}
Define positive functions recursively by
\begin{equation}\label{path:chains}
 \ell_{i+1}(a')=\int\ell_i(a)\mathcal K(a,a')\,d\nu(a),\qquad
 h_{k+1}(a)=\int\mathcal K(a,a')h_k(a')\,d\nu(a').
\end{equation}
All integrals without limits in this section are over $[-1,1]$. If $Z_j$ is
the volume of the moment polytope, then
\[
 Z_j=\int\ell_i(a)h_{j-1-i}(a)\,d\nu(a),\qquad D=n!Z_j.
\]
For a hexagon ray $v$ at position $i$, and an interior base ray at position $c$,
respectively, the normalized degrees are
\begin{align}
 \delta_{i,v}&=\frac{\int\ell_i h_{j-1-i}\,dE_v}
                         {\int\ell_i h_{j-1-i}\,d\nu},\label{path:hexdegree}\\
 \delta_c&=\frac{\iint\ell_{c-1}(a)(3+a-a')h_{j-1-c}(a')\,d\nu(a)d\nu(a')}
 {\iint\ell_{c-1}(a)\mathcal K(a,a')h_{j-1-c}(a')\,d\nu(a)d\nu(a')}.
 \label{path:basedegree}
\end{align}
Here $E_v$ is the push-forward of lattice length on the corresponding hexagon
edge. For $v=r,-r$ it is unit mass at $-1,1$, respectively; for the other four
rays it is Lebesgue measure on $[-1,0]$ or $[0,1]$, twice each.
At the left end, the numerator for a single ray is obtained by replacing
$\ell_0$ with $(b_0+1-a)^{b_0-1}/(b_0-1)!$; the right end is analogous.
These formulas follow by integrating the product of simplex volumes in
Section~\ref{sec:graphs}.

For completeness, the rational matrices needed to evaluate the integrals are
as follows. On $1,a,a^2$, the moment matrix and the kernel coefficient matrix are
\begin{equation}\label{path:matrices}
 W=\begin{pmatrix}3&0&5/6\\0&5/6&0\\5/6&0&7/15\end{pmatrix},\qquad
 \kappa=\begin{pmatrix}9/2&-3&1/2\\3&-1&0\\1/2&0&0\end{pmatrix}.
\end{equation}
For cubic end factors, append a fourth row and column, using
\[
 W_{pq}=I_{p+q},\qquad
 I_t=\begin{cases}0&t\text{ odd},\\
 2\bigl(2/(t+1)-1/(t+2)\bigr)&t\text{ even},\end{cases}
 \qquad 0\le p,q\le3,
\]
and extend $\kappa$ by zero. If $\lambda,\eta$ are the coefficient vectors of
$\ell_0,h_0$, then
\begin{equation}\label{path:matrixvolume}
 Z_j=\lambda^{\mathsf T}(W\kappa)^{j-1}W\eta.
\end{equation}
Replace the relevant $W$ by $W_v=(\int a^{p+q}\,dE_v)_{pq}$ to compute a hexagon
degree. For an interior base degree replace the relevant $\kappa$ by the
coefficient matrix of $3+a-a'$. The four edge moment sequences are
$(-1)^t$, $1$, $(-1)^t/(t+1)$, and $1/(t+1)$. Thus every finite-path degree
used below is rational and is determined by the displayed data.

\subsection{Invariant subspaces and ranks}
Let $A$ be generated by the independent hexagon reflections
$(x,y)\mapsto(x-y,-y)$, the permutations of the two untwisted rays at each
interior base factor, and the permutations of the $b$ untwisted rays at each
end $\PP^b$. These preserve the fan and the degrees. By
Proposition~\ref{prop:symmetry}, it suffices to test the $A$-invariant
ray-spanned subspaces, since the path has connected ray matroid.

Taking all rays in the span of a ray subset preserves its rank and increases
its weight. The two nontrivial reflection orbits in a hexagon each span its
plane. Hence it suffices to consider the following \emph{configurations}:
choose, at each hexagon, nothing, the pair $\{r_i,-r_i\}$, or all six rays;
at each base factor, choose all its untwisted rays or none, and choose its
twisted ray independently. Configurations need not be closed under taking
the rays in their span, but include every invariant closed set. Positivity
for all nonempty configurations of proper rank is therefore equivalent to stability.

Write $H_i\in\{0,1,2\}$ for the hexagon choice, and $(C_c,T_c)\in\{0,1\}^2$
for the choices of the untwisted rays and twisted ray at base vertex $c$.
Let $g_c=\sum_{t=1}^{b_c}f_t^c$. The nontrivial representation contributed by
a selected hexagon plane has dimension one; that contributed by selected
untwisted base rays has dimension $b_c-1$. For $b_c=1$ this latter
contribution is zero. The remaining vectors lie in the space with basis
$r_0,\ldots,r_{j-1},g_0,\ldots,g_j$, where
\begin{equation}\label{path:relations}
 f_0^0=-g_0-r_0,\qquad
 f_0^c=-g_c+r_{c-1}-r_c\ (1\le c<j),\qquad
 f_0^j=-g_j+r_{j-1}.
\end{equation}

\begin{lemma}\label{path:rank}
The rank of a configuration is computed as follows. A selected hexagon line
or plane contributes $H_i$; selected untwisted rays at vertex $c$ contribute
$b_c C_c$. Each selected twisted ray with $C_c=0$ contributes one further
dimension. For the remaining twisted rays, form a graph with vertices
$r_0,\ldots,r_{j-1}$ and one additional vertex $*$. Identify every selected
$r_i$ with $*$. Include $r_{c-1}r_c$ when $C_cT_c=1$, and include the edges
$*r_0$ and $r_{j-1}*$ when $C_0T_0=1$ and $C_jT_j=1$. Add the rank of this
graphic matroid to the preceding contributions.

A configuration spans the whole lattice space if and only if every hexagon
is a plane, every interior base has $C_c=1$, and each end has $C_c=1$ or
has $b_c=1$ and $T_c=1$.
\end{lemma}

\begin{proof}
The nontrivial representations of the independent group factors are distinct
and occur with multiplicity one. After accounting for them and the selected
basis vectors $r_i,g_c$, a twisted ray whose $g_c$ is not selected has a
coordinate occurring in no other remaining vector. It adds one dimension
and may be removed. The other vectors in \eqref{path:relations}, modulo
the selected basis vectors, are precisely the incidence vectors of the
stated graph. Their rank is its graphic-matroid rank. Full rank requires
each nontrivial representation, hence every hexagon plane and every
interior untwisted pair. All $r_i$ are then selected, and the only remaining
condition is to obtain each end's base space, as stated.
\end{proof}

Here is an explicit recurrence for the minimum. After processing hexagon $i$,
the state of $r_i$ is $G$ if it is selected, $g$ if it is unselected but joined
to $*$ through the processed edges, and $u$ otherwise. A period consists of
an interior base factor followed by its right-hand hexagon. Let $w_h$ be the
weight of hexagon choice $h$, where $w_0=0$, and let $b$ be the weight of one
base ray. For a choice $(C,T,h)$ put
\[
 L=CT,\qquad a(C,T)=2C+T(1-C).
\]
The rank increment, new state, and weight increment are given by
\begin{equation}\label{path:transition}
\begin{array}{c|c|c|c}
\text{condition}&\text{rank increment}&\text{new state}&\text{weight increment}\\ \hline
h\ge1&a(C,T)+h+L\,\mathbf1_{s=u}&G&(2C+T)b+w_h\\
h=0,\ L=1&a(C,T)+1&g\text{ if }s\ne u,\ u\text{ otherwise}&(2C+T)b\\
h=0,\ L=0&a(C,T)&u&(2C+T)b
\end{array}
\end{equation}
The left end contributes $b_0C_0+T_0(1-C_0)$ dimensions before the first
hexagon; its possible remaining edge starts at $*$. The right end contributes
$b_jC_j+T_j(1-C_j)+C_jT_j\mathbf1_{s=u}$ dimensions.
Its weight is $(b_jC_j+T_j)\delta_j$, and similarly at the left. To process
the initial hexagon, use \eqref{path:transition} with incoming state $G$,
$L=C_0T_0$, and $a(C,T)=0$, since the left base contribution has already
been included; only the hexagon weight is subtracted at this step.

Keep Boolean variables $e,f$ recording whether any ray has been selected
and whether the full-rank conditions on the processed factors still hold.
After choosing the left end, initialize
\[
 e=C_0\lor T_0,\qquad
 f=C_0\lor(b_0=1\land T_0).
\]
After the initial hexagon, replace these by $e\lor(H_0>0)$ and
$f\land(H_0=2)$. Each interior period updates them by
\[
 e\longmapsto e\lor C\lor T\lor(h>0),\qquad
 f\longmapsto f\land C\land(h=2).
\]
At the right end, the final values are $e\lor C_j\lor T_j$ and
$f\land\bigl(C_j\lor(b_j=1\land T_j)\bigr)$. At each step retain the
least cost, defined as rank minus weight, for each state and these two
variables. At the end discard final $e=0$ and final $f=1$. This gives the
exact \emph{minimum invariant rank deficit}, denoted $m_A(j;b_0,b_j)$.
Only its positivity is needed; no equality with the minimum over arbitrary
subspaces is asserted.

\subsection{Uniform convergence}
For positive functions on $[-1,1]$ define Hilbert's projective distance by
\[
 d_H(f,g)=\log\left(\sup_a\frac{f(a)}{g(a)}\sup_a\frac{g(a)}{f(a)}\right),
 \qquad \Delta=\log(81/25),\quad \tau=2/7.
\]

\begin{lemma}\label{path:contraction}
For $i,i'\ge1$ one has
$d_H(\ell_i,\ell_{i'})\le\Delta\tau^{\min(i,i')-1}$, and the same holds for
the right-hand functions. The bound also compares chains starting from
different end factors. For a hexagon or interior-base degree ratio, viewed
as a function $\delta(i,k)$ of its left and right chain indices,
\[
 |\log\delta(i,k)-\log\delta(i',k')|
 \le d_H(\ell_i,\ell_{i'})+d_H(h_k,h_{k'}).
\]
The limits $\delta(i,\infty)$, $\delta(\infty,k)$, and
$\delta_\infty=\delta(\infty,\infty)$ exist. In particular,
\begin{align*}
 |\log\delta(i,k)-\log\delta_\infty|
   &\le\Delta(\tau^{i-1}+\tau^{k-1}) &&(i,k\ge1),\\
 |\log\delta(i,k)-\log\delta(i,\infty)|
   &\le\Delta\tau^{k-1} &&(i\ge0,\ k\ge1).
\end{align*}
These bulk limits are independent of the end factors. An end-ray degree
has a one-sided limit, obtained by sending the opposite chain index to
infinity; its logarithmic error satisfies the corresponding one-index
bound above, with the end factor fixed. This limit may depend on that end
factor. Every individual ray weight, including all these limiting
weights, is at most $27$.
\end{lemma}

\begin{proof}
The Birkhoff--Hopf theorem for a strictly positive continuous kernel gives
projective diameter equal to the logarithm of its maximal cross-ratio and
contraction factor $\tanh(\Delta/4)$
\cite[Theorems 3.5 and 6.3]{EvesonNussbaum1995}.
The measure $\nu$ has full support. For $a>b$, the cross-ratio of
$3+a-a'$ increases with $a'$ and decreases with $b'$; at $a'=1,b'=-1$ it is
\[
 \frac{(2+a)(4+b)}{(4+a)(2+b)}\le\frac95.
\]
This increases with $a$ and decreases with $b$, and equality holds at
$a=1,b=-1$. Exchanging the pairs treats $a<b$. Squaring gives $81/25$
for $\mathcal K$, and $\tanh(\Delta/4)=2/7$. One application puts any two
positive initial functions at distance at most $\Delta$; each subsequent
application contracts this distance by $\tau$.

Each degree ratio is the ratio of two integrals of the chain functions
against positive measures. If $m\le f/\widetilde f\le M$, replacing $f$
by $\widetilde f$ changes such a ratio by a factor between $m/M$ and $M/m$.
Applying this to the two chain functions proves the logarithmic estimate.
It makes their degree ratios Cauchy and proves the stated limits without
any spectral assumption.

For the last assertion, each initial end polynomial has maximum divided by
minimum at most $9$. For fixed $a$, the ratio of the largest to the smallest
value of $\mathcal K(a,a')$ as $a'$ varies is at most $9$; the transposed
kernel has the same property. All chain functions therefore have oscillation
ratio at most $9$. A hexagon edge measure has total mass $1$ and $\nu$ has
mass $3$, so \eqref{path:hexdegree} is at most $9^2/3=27$. An interior
base weight is an average of $2/(3+a-a')\le2$, and an end weight is an
average of $b/(b+1\pm a)\le1$. The bounds persist under limits.
\end{proof}

\subsection{The finite calculation and its uniform error bound}
At limiting weights an interior period has total weight exactly $4$.
Indeed, Lemma~\ref{path:contraction} bounds the difference between the total
weight of a path and $(j-1)$ times its limiting period weight by a constant
independent of $j$: the geometric tails are summable and the end factors
have bounded weight. The total weight is $n=4j-2+b_0+b_j$, so division by
$j$ gives the assertion.

Consequently the all-selected period at state $G$ has cost zero, as does
the empty period at state $u$. For the limiting period recurrence, every
other simple directed cycle has cost greater than $0.1259$. This is a
finite rational verification with controlled error: substitute the chain
indices $(80,80)$ for $(\infty,\infty)$ in
\eqref{path:hexdegree}--\eqref{path:basedegree}, enumerate the $36$
transitions of \eqref{path:transition}, and then the simple cycles of
length at most three. Table~\ref{path:cycles} lists the least transition
costs after removing the two zero loops. Each positive entry lies within
$10^{-6}$ of the displayed number; the limiting entries marked $0$ are
exactly zero. There is no transition from $u$ to $g$. Thus those remaining
zero transitions lie on no zero-cost cycle.

\begin{table}[ht]
\centering
\caption{Least limiting transition costs, with the all-selected $G\to G$
and empty $u\to u$ loops removed.}
\label{path:cycles}
\begin{adjustbox}{max width=\textwidth,center}
\begin{tabular}{c|ccc}
&$G$&$g$&$u$\\\hline
$G$&0.125902&0.874098&0\\
$g$&0&0.874098&0\\
$u$&0.125902&---&0.291366
\end{tabular}
\end{adjustbox}
\end{table}

We now specify the finite comparison for a path of length $j=2q+m$, where
$q=40$ and $m\ge0$. A period has index $p$ when it consists of the base
factor $p$ and hexagon $p$. Define comparison weights by using
$\delta(i,\infty)$ at the left end and periods $1,\ldots,q-1$,
$\delta_\infty$ at periods $q,\ldots,q+m-1$, and
$\delta(\infty,k)$ at periods $q+m,\ldots,j-1$ and the right end.
The initial hexagon belongs to the left end. Let $M(m)$ denote the costs
of configurations with these comparison weights.

There are at most $9q+4$ rays in each end zone. At a hexagon $p$ the chain
indices are $(p,j-1-p)$; at an interior base factor $p$ they are
$(p-1,j-1-p)$. Put
\[
 \overline\Delta=\frac{11757}{10000}>\Delta,\quad
 x=\overline\Delta\tau^{q-2},\quad
 y=2\overline\Delta\tau^{79},\quad s=\frac{2x}{1-\tau}.
\]
The strict upper bound on $\Delta$ follows, using rational arithmetic, from
$\sum_{r=0}^7\overline\Delta^r/r!>81/25$.
Since $e^t-1\le t+t^2$ for $0\le t\le1$, Lemma~\ref{path:contraction}
bounds the sum of all weight changes between a true path and $M(m)$ by
\begin{equation}\label{path:error}
 E_{\rm path}=27\{9s(1+2x)+2(9q+4)(x+x^2)\}.
\end{equation}
The first term sums the two geometric tails in the central periods. The
second treats both end zones, using the weaker index $q-2$ that also
applies to the first base factor of the right zone when $m=0$.
This bound is uniform in $m$ and in the configuration selected.

For the four comparison paths $M(0),\ldots,M(3)$, replacing every infinite
chain index by $80$ changes the cost of any configuration by at most
\begin{equation}\label{path:proxyerror}
 E_{\rm rat}=27\{9(2q+4)+4\}(y+y^2)<6\cdot10^{-39}.
\end{equation}
The same bound more than covers a cycle of length at most three. All such
substituted computations use rational numbers only. Their all-selected
costs have absolute value below $10^{-36}$. Thus the true all-selected
cost of $M(m)$ has absolute value below $2\cdot10^{-36}$ for every $m$,
because adding a limiting full period changes that cost by zero.

Table~\ref{path:bounds} gives rational lower bounds from the recurrence,
first for every true path $1\le j\le84$, and then for each of the four
rational comparison paths. The bounds in the table are terminating decimals,
not rounded approximations to the minima. Reversing the path exchanges its
ends, so the six rows cover all nine ordered end types.

\begin{table}[ht]
\centering
\caption{Lower bounds for the minimum invariant rank deficit of nonempty
configurations of proper rank. The comparison column uses chain index $80$
in place of infinity.}
\label{path:bounds}
\begin{adjustbox}{max width=\textwidth,center}
\begin{tabular}{ccc}
\toprule
$(b_0,b_j)$& every true path $1\le j\le84$ & each comparison path $M(0),\ldots,M(3)$\\
\midrule
\csname @@input\endcsname tables/path_bounds.tex
\bottomrule
\end{tabular}
\end{adjustbox}
\end{table}

The calculation is completely specified by \eqref{path:matrices}, the moment
sequences following it, the end polynomials \eqref{path:kernel}, and the
recurrence \eqref{path:transition}, with the two rank flags described there.
The accompanying implementation reproduces these tables using exact rational
arithmetic; Appendix~\ref{sec:computation} describes additional independent
checks.

\begin{proof}[Proof of Theorem~\ref{path:stable}]
The finite calculation proves positivity for $j\le84$. For the remaining
lengths use $M(m)$ as above. Every configuration of proper nonzero rank in
$M(0),\ldots,M(3)$ costs more than $0.0295-E_{\rm rat}$. The empty
configuration costs zero. A full-rank configuration costs at least the
all-selected configuration, because all weights are positive, and hence
at least $-2\cdot10^{-36}$.

For $m\ge4$, two boundaries of the central periods have the same state.
Delete the periods between them. The remaining configuration belongs to
$M(m')$ with $m'<m$ and has the same rank increments outside the deleted
walk. If the walk consists only of the all-selected zero loop or only of
the empty zero loop, deletion preserves both its cost and whether the
whole configuration is nonempty of proper rank. In the first case all
deleted hexagons and base factors satisfy the full-rank conditions; the
remaining state $G$ prevents the new configuration from being empty. In
the second case the remaining state $u$ prevents it from having full rank.
If the walk is not of either kind, its decomposition into simple cycles
contains a positive-cost cycle and its cost exceeds $0.1259$.

Induction therefore gives a lower bound $0.0295-E_{\rm rat}-2\cdot10^{-36}$
for every nonempty proper-rank configuration of $M(m)$. Subtracting the
true-path error, every relevant rank deficit is therefore greater than
\[
 0.0295-\bigl(E_{\rm path}+2E_{\rm rat}+2\cdot10^{-36}\bigr)
 >0.0295-6\cdot10^{-17}>0.
\]
Apply
Proposition~\ref{prop:symmetry} to conclude stability.
\end{proof}

\begin{corollary}\label{path:everydimension}
For every $n\ge4$ there is a smooth toric Fano $n$-fold with anticanonically
stable tangent bundle and Picard number $5\lfloor n/4\rfloor+1$.
This is the largest Picard number of a stable example among the connected
Fano hexagon bundles over graphs defined in Section~\ref{sec:graphs}.
\end{corollary}

\begin{proof}
Write $n=4j+t$, $0\le t\le3$. Choose the end types
$(1,1),(2,1),(2,2),(3,2)$ for $t=0,1,2,3$, respectively.
Equation~\eqref{path:dimension} gives the claimed dimension and Picard
number, and Theorem~\ref{path:stable} gives stability. The upper bound
within this class is the graph bound of Section~\ref{sec:graphs}.
\end{proof}
